% segment-local-id: OR03-S008
% segment-id: d90.orig.v1.tr03.seg0008
% license: CC BY-SA 4.0
% provenance: Materi asesmen asli berbahasa Indonesia, disusun secara independen.

\section{Set soal VI: transportasi optimal dan sintesis}
\label{orig03:ps06}

% stable-id: d90.orig.v1.tr03.ps06.exercise.0001
\begin{exercise}[Stabilitas terhadap perturbasi biaya seragam]
\label{orig03:ps06:exercise:0001}
Misalkan $\Pi(\alpha,\beta)$ adalah himpunan kopling probabilitas dan dua
fungsi biaya memenuhi
$\lVert c-\widetilde c\rVert_\infty\le\varepsilon$. Tuliskan $v$ dan
$\widetilde v$ untuk nilai optimal dengan biaya $c$ dan $\widetilde c$.

\begin{enumerate}
  \item Buktikan $|v-\widetilde v|\le\varepsilon$.
  \item Jika $P_c$ optimal untuk biaya $c$, buktikan bahwa ia paling jauh
  $2\varepsilon$-suboptimal untuk biaya $\widetilde c$.
\end{enumerate}
\end{exercise}

% stable-id: d90.orig.v1.tr03.ps06.hint.0001
\paragraph{Petunjuk bertahap.}
\label{orig03:ps06:hint:0001}
\textbf{Tahap 1.} Integralkan batas titik demi titik terhadap sembarang kopling
probabilitas.
\textbf{Tahap 2.} Sisipkan nilai $v$ di antara biaya
$\widetilde c$ pada $P_c$ dan nilai $\widetilde v$.

% stable-id: d90.orig.v1.tr03.ps06.answer.0001
\paragraph{Jawaban singkat.}
\label{orig03:ps06:answer:0001}
Untuk setiap $P\in\Pi(\alpha,\beta)$,
$|\langle c,P\rangle-\langle\widetilde c,P\rangle|\le\varepsilon$.
Mengambil infimum memberi $|v-\widetilde v|\le\varepsilon$, dan
\[
 \langle\widetilde c,P_c\rangle-\widetilde v\le2\varepsilon.
\]

% stable-id: d90.orig.v1.tr03.ps06.solution.0001
\paragraph{Solusi lengkap.}
\label{orig03:ps06:solution:0001}
Karena massa total setiap kopling adalah satu,
\[
 \left|\int(c-\widetilde c)\,dP\right|
 \le\int|c-\widetilde c|\,dP\le\varepsilon.
\]
Jadi $\langle\widetilde c,P\rangle\le\langle c,P\rangle+\varepsilon$
untuk setiap $P$. Mengambil infimum memberi
$\widetilde v\le v+\varepsilon$. Menukar peran kedua biaya memberi
$v\le\widetilde v+\varepsilon$, sehingga batas nilai optimal terbukti.
Selanjutnya,
\[
 \langle\widetilde c,P_c\rangle-\widetilde v
 \le\langle c,P_c\rangle+\varepsilon-(v-\varepsilon)
 =2\varepsilon,
\]
karena $\langle c,P_c\rangle=v$. Batas kedua memuat dua sumber galat: evaluasi
rencana dan pergeseran nilai optimal.

% stable-id: d90.orig.v1.tr03.ps06.exercise.0002
\begin{exercise}[Sertifikat dual untuk translasi]
\label{orig03:ps06:exercise:0002}
Pada $\mathbb R$, ambil biaya
$c(x,y)=\tfrac12(y-x)^2$. Misalkan $\alpha$ mempunyai momen pertama hingga dan
$\beta=(x\mapsto x+a)_\#\alpha$ untuk suatu $a\in\mathbb R$.

\begin{enumerate}
  \item Hitung biaya kopling deterministik $y=x+a$.
  \item Buktikan bahwa
  $\phi(x)=-ax$ dan $\psi(y)=ay-a^2/2$ layak untuk masalah dual Kantorovich.
  \item Gunakan kesetaraan primal--dual untuk menyimpulkan optimalitas kopling
  translasi tersebut.
\end{enumerate}
\end{exercise}

% stable-id: d90.orig.v1.tr03.ps06.hint.0002
\paragraph{Petunjuk bertahap.}
\label{orig03:ps06:hint:0002}
\textbf{Tahap 1.} Kurangi $\phi(x)+\psi(y)$ dari $c(x,y)$ dan lengkapi
kuadrat dalam $y-x$.
\textbf{Tahap 2.} Gunakan identitas
$\int y\,d\beta(y)=\int(x+a)\,d\alpha(x)$.

% stable-id: d90.orig.v1.tr03.ps06.answer.0002
\paragraph{Jawaban singkat.}
\label{orig03:ps06:answer:0002}
Biaya kopling translasi adalah $a^2/2$, dan
\[
 c(x,y)-\phi(x)-\psi(y)=\tfrac12(y-x-a)^2\ge0.
\]
Nilai dual kedua potensial juga $a^2/2$, sehingga kopling dan potensial itu
optimal.

% stable-id: d90.orig.v1.tr03.ps06.solution.0002
\paragraph{Solusi lengkap.}
\label{orig03:ps06:solution:0002}
Pada kopling $y=x+a$, biaya bernilai konstan
$\tfrac12a^2$, sehingga integralnya $a^2/2$. Untuk semua $x,y$,
\[
 \tfrac12(y-x)^2-\bigl(-ax+ay-a^2/2\bigr)
 =\tfrac12(y-x-a)^2\ge0.
\]
Maka $\phi(x)+\psi(y)\le c(x,y)$, dengan kesetaraan tepat pada graf translasi.
Momen pertama menjamin integral potensial terdefinisi, dan
\[
 \int\phi\,d\alpha+\int\psi\,d\beta
 =-a\int x\,d\alpha
  +a\int(x+a)\,d\alpha-a^2/2
 =a^2/2.
\]
Dual lemah menyatakan bahwa nilai ini tidak dapat melebihi biaya kopling mana
pun. Karena kopling translasi mencapai nilai yang sama, celah dualnya nol dan
kopling itu optimal.

% stable-id: d90.orig.v1.tr03.ps06.exercise.0003
\begin{exercise}[Penskalaan Sinkhorn dalam domain log dan kebebasan gauge]
\label{orig03:ps06:exercise:0003}
Untuk marginal positif $a\in\mathbb R^n$ dan $b\in\mathbb R^m$, definisikan
\[
 K_{ij}=\exp(-C_{ij}/\varepsilon),
 \qquad P=\operatorname{diag}(u)K\operatorname{diag}(v),
 \]
dengan $u_i,v_j>0$. Tuliskan $\rho_i=\log u_i$ dan $\tau_j=\log v_j$.

\begin{enumerate}
  \item Turunkan pembaruan baris dan kolom dalam bentuk log-sum-exp.
  \item Buktikan bahwa $u\mapsto tu$, $v\mapsto v/t$ tidak mengubah $P$.
  \item Jelaskan secara matematis mengapa pemusatan $\rho$ dan $\tau$ dengan
  pergeseran yang berlawanan dapat mencegah luapan numerik tanpa mengubah
  kopling.
\end{enumerate}
\end{exercise}

% stable-id: d90.orig.v1.tr03.ps06.hint.0003
\paragraph{Petunjuk bertahap.}
\label{orig03:ps06:hint:0003}
\textbf{Tahap 1.} Selesaikan persamaan $P\1=a$ terhadap setiap $u_i$,
lalu ambil logaritma; lakukan hal yang sama untuk kolom.
\textbf{Tahap 2.} Tambahkan konstanta $s$ pada semua $\rho_i$ dan kurangkan
$s$ dari semua $\tau_j$.

% stable-id: d90.orig.v1.tr03.ps06.answer.0003
\paragraph{Jawaban singkat.}
\label{orig03:ps06:answer:0003}
\[
 \rho_i=\log a_i-\operatorname{LSE}_j(-C_{ij}/\varepsilon+\tau_j),
 \quad
 \tau_j=\log b_j-\operatorname{LSE}_i(-C_{ij}/\varepsilon+\rho_i).
\]
Transformasi $(\rho,\tau)\mapsto(\rho+s\1,\tau-s\1)$ membiarkan
setiap jumlah $\rho_i+\tau_j$ tetap, sehingga $P$ tidak berubah.

% stable-id: d90.orig.v1.tr03.ps06.solution.0003
\paragraph{Solusi lengkap.}
\label{orig03:ps06:solution:0003}
Kendala baris adalah
$u_i\sum_jK_{ij}v_j=a_i$, jadi
\[
 \rho_i=\log a_i-
 \log\sum_j\exp(-C_{ij}/\varepsilon+\tau_j).
\]
Logaritma jumlah eksponensial adalah fungsi $\operatorname{LSE}$. Dengan cara
yang sama, kendala kolom menghasilkan rumus untuk $\tau_j$ pada jawaban
singkat. Jika $t>0$, maka
\[
 \operatorname{diag}(tu)K\operatorname{diag}(v/t)
 =\operatorname{diag}(u)K\operatorname{diag}(v).
\]
Dalam koordinat log, ini tepat sama dengan menambahkan $s=\log t$ pada
$\rho$ dan mengurangkannya dari $\tau$. Karena entri $P$ bergantung pada
$\rho_i+\tau_j-C_{ij}/\varepsilon$, semua entri tetap sama. Kita boleh memilih
$s$, misalnya untuk memusatkan rentang nilai log, agar eksponensial perantara
tidak terlalu besar atau kecil. Operasi itu mengubah representasi, bukan
kopling atau marginalnya.

% stable-id: d90.orig.v1.tr03.ps06.exercise.0004
\begin{exercise}[Batas bias regularisasi entropi]
\label{orig03:ps06:exercise:0004}
Pada tabel kopling probabilitas berukuran $n\times m$, definisikan
$H(P)=-\sum_{ij}P_{ij}\log P_{ij}$ dengan konvensi $0\log0=0$.
Misalkan $P^\star$ meminimalkan $\langle C,P\rangle$ dan $P_\varepsilon$
meminimalkan
\[
 \langle C,P\rangle-\varepsilon H(P)
\]
pada himpunan kopling yang sama. Buktikan
\[
 0\le\langle C,P_\varepsilon\rangle-
       \langle C,P^\star\rangle
 \le\varepsilon\log(nm).
\]
\end{exercise}

% stable-id: d90.orig.v1.tr03.ps06.hint.0004
\paragraph{Petunjuk bertahap.}
\label{orig03:ps06:hint:0004}
\textbf{Tahap 1.} Bandingkan nilai objektif teratur pada $P_\varepsilon$ dan
$P^\star$.
\textbf{Tahap 2.} Gunakan $0\le H(P)\le\log(nm)$ untuk setiap distribusi pada
$nm$ sel.

% stable-id: d90.orig.v1.tr03.ps06.answer.0004
\paragraph{Jawaban singkat.}
\label{orig03:ps06:answer:0004}
Optimalitas kedua rencana memberi batas bawah nol dan
\[
 \langle C,P_\varepsilon\rangle-langle C,P^\star\rangle
 \le\varepsilon\bigl(H(P_\varepsilon)-H(P^\star)\bigr)
 \le\varepsilon\log(nm).
\]

% stable-id: d90.orig.v1.tr03.ps06.solution.0004
\paragraph{Solusi lengkap.}
\label{orig03:ps06:solution:0004}
Karena $P^\star$ optimal untuk biaya tanpa regularisasi,
$\langle C,P_\varepsilon\rangle\ge\langle C,P^\star\rangle$. Di sisi lain,
optimalitas $P_\varepsilon$ untuk objektif teratur menyatakan
\[
 \langle C,P_\varepsilon\rangle-\varepsilon H(P_\varepsilon)
 \le\langle C,P^\star\rangle-\varepsilon H(P^\star).
\]
Menata ulang memberi ketaksamaan pertama pada jawaban singkat. Entropi
distribusi dengan paling banyak $nm$ atom berada di antara nol dan
$\log(nm)$. Maka
$H(P_\varepsilon)-H(P^\star)\le\log(nm)$, yang menyelesaikan batas atas.
Batas ini mengukur bias biaya tak teratur, bukan jarak antara kedua rencana.

% stable-id: d90.orig.v1.tr03.ps06.exercise.0005
\begin{exercise}[Kelengkungan dual entropik dan arah gauge]
\label{orig03:ps06:exercise:0005}
Pertimbangkan fungsi dual
\[
 D(f,g)=a^\top f+b^\top g
 -\varepsilon\sum_{i,j}
 \exp\!\left(\frac{f_i+g_j-C_{ij}}{\varepsilon}\right),
\]
dan tuliskan
$P_{ij}=\exp((f_i+g_j-C_{ij})/\varepsilon)$.

\begin{enumerate}
  \item Hitung gradien terhadap $f$ dan $g$.
  \item Hitung bentuk kuadratik Hessian pada arah $(s,t)$ dan buktikan
  kekonkafan.
  \item Identifikasi arah gauge. Jika graf bipartit dukungan positif $P$
  terhubung, buktikan bahwa hanya arah gauge yang mempunyai kelengkungan nol.
\end{enumerate}
\end{exercise}

% stable-id: d90.orig.v1.tr03.ps06.hint.0005
\paragraph{Petunjuk bertahap.}
\label{orig03:ps06:hint:0005}
\textbf{Tahap 1.} Diferensiasikan satu eksponensial dan kelompokkan jumlah
baris serta kolom.
\textbf{Tahap 2.} Bentuk kuadratik dapat ditulis sebagai jumlah berbobot dari
$(s_i+t_j)^2$.

% stable-id: d90.orig.v1.tr03.ps06.answer.0005
\paragraph{Jawaban singkat.}
\label{orig03:ps06:answer:0005}
\[
 \nabla_fD=a-P\1,
 \qquad \nabla_gD=b-P^\top\1,
\]
dan
\[
 \begin{pmatrix}s\\t\end{pmatrix}^{\!\top}
 \nabla^2D
 \begin{pmatrix}s\\t\end{pmatrix}
 =-\frac1\varepsilon\sum_{i,j}P_{ij}(s_i+t_j)^2\le0.
\]
Arah $(s,t)=(c\1,-c\1)$ adalah arah gauge; bila dukungan
terhubung, semua arah nol berbentuk demikian.

% stable-id: d90.orig.v1.tr03.ps06.solution.0005
\paragraph{Solusi lengkap.}
\label{orig03:ps06:solution:0005}
Diferensiasi terhadap $f_i$ menghasilkan
$a_i-\sum_jP_{ij}$, sedangkan diferensiasi terhadap $g_j$ menghasilkan
$b_j-\sum_iP_{ij}$. Diferensiasi kedua pada lintasan
$(f+\theta s,g+\theta t)$ memberi
\[
 \left.\frac{d^2}{d\theta^2}D(f+\theta s,g+\theta t)\right|_{\theta=0}
 =-\frac1\varepsilon\sum_{i,j}P_{ij}(s_i+t_j)^2.
\]
Bobotnya nonnegatif, sehingga $D$ konkaf. Bentuk kuadratik nol tepat bila
$s_i+t_j=0$ pada setiap sisi dukungan positif. Pada graf bipartit yang
terhubung, kondisi ini memaksa semua $s_i$ sama dengan satu konstanta $c$ dan
semua $t_j=-c$. Itulah transformasi
$(f,g)\mapsto(f+c\1,g-c\1)$ yang tidak mengubah $P$. Jadi dual
bersifat konkaf ketat setelah satu gauge dipatok, meskipun tidak konkaf ketat
pada representasi yang belum dinormalisasi.
